\documentclass[11pt]{report}

% =====================================================================
% PREAMBLE
% =====================================================================

\usepackage[T1]{fontenc}
\usepackage[utf8]{inputenc}
\usepackage[english]{babel}

\usepackage{graphicx}
\usepackage{geometry}
\usepackage{amsmath}

% Code listings
\usepackage{listings}
\usepackage{xcolor}

\usepackage{mdframed}
\usepackage{adjustbox}
\usepackage{hyperref}

% Page geometry
\geometry{
    a4paper,
    top=2.2cm,
    bottom=2.2cm,
    left=2.4cm,
    right=2.4cm
}

\usepackage{titlesec}
\usepackage{enumitem}
\usepackage{parskip}

\titlespacing*{\chapter}{0pt}{-20pt}{10pt}
\titlespacing*{\section}{0pt}{10pt}{4pt}
\titlespacing*{\subsection}{0pt}{8pt}{3pt}
\titlespacing*{\subsubsection}{0pt}{6pt}{2pt}

\setlist{itemsep=1pt, topsep=2pt, parsep=0pt}
\usepackage[T1]{fontenc}
\usepackage[utf8]{inputenc}
\usepackage[english]{babel}
\usepackage{graphicx}
\usepackage{geometry}
\usepackage{hyperref}
\geometry{a4paper, top=3cm, bottom=3cm, left=3cm, right=3cm}

\hypersetup{
    colorlinks=true,
    linkcolor=black,
    filecolor=black,
    urlcolor=blue,
    citecolor=black,
}

\usepackage[expansion=false]{microtype}
\usepackage{fancyvrb}
\usepackage{fvextra}

\setlength{\parskip}{4pt}
\setlength{\parindent}{0pt}


% =====================================================================
% CODE LISTING STYLE
% =====================================================================

\definecolor{codebackground}{RGB}{245,245,245}

\lstdefinestyle{cppstyle}{
    language=C++,
    backgroundcolor=\color{codebackground},
    basicstyle=\ttfamily\footnotesize,
    keywordstyle=\bfseries,
    commentstyle=\itshape,
    numbers=left,
    numberstyle=\tiny,
    numbersep=8pt,
    frame=single,
    framesep=4pt,
    breaklines=true,
    tabsize=4,
    showstringspaces=false,
    columns=fullflexible,
    aboveskip=4pt,
    belowskip=4pt
}

\lstset{style=cppstyle}

\lstdefinelanguage{CSharp}{
    language=[Sharp]C,
    morekeywords={var, async, await, yield, get, set, value, nameof}
}

% Clean, academic style for all code blocks
\lstset{
    basicstyle=\ttfamily\footnotesize,
    keywordstyle=\bfseries\color{blue!70!black},
    commentstyle=\itshape\color{green!50!black},
    stringstyle=\color{red!60!black},
    showstringspaces=false,
    numbers=none, % set to 'left' if line numbers are desired
    frame=single, % frame around the code
    breaklines=true,
    aboveskip=4pt,
    belowskip=4pt
}

% =====================================================================
% VERBATIM STYLE FOR PROMPTS (uniform, margin-safe)
% =====================================================================

% Ensures every \begin{Verbatim} block used for prompts wraps long lines
% instead of overflowing the page margins, and keeps a consistent visual
% style (small monospaced font, light background, thin frame) across the
% whole document.
\DefineVerbatimEnvironment{PromptVerbatim}{Verbatim}{
    breaklines=true,
    breakanywhere=true,
    breaksymbolleft=,
    fontsize=\small,
    frame=single,
    framesep=3mm,
    xleftmargin=0pt,
    xrightmargin=0pt,
    samepage=false
}

% =====================================================================
% EXAMPLE BOX STYLE
% =====================================================================

\newmdenv[
    leftline=true,
    rightline=false,
    topline=false,
    bottomline=false,
    linecolor=black,
    linewidth=2pt,
    innerleftmargin=12pt,
    innerrightmargin=8pt,
    innertopmargin=8pt,
    innerbottommargin=8pt
]{examplebox}


% =====================================================================
% MAIN DOCUMENT
% =====================================================================

\begin{document}
\input{title_page/Frontespizio_prompt_log} % include the translated title page

\chapter{AI Prompt Engineering}
\label{AIPromptEngineering}

To ensure the technical accuracy of the generated content and its adherence
to the project guidelines, this report was developed using a
\textbf{two-phase prompting strategy}: prompt improvement and prompt
execution. A fundamental principle of this methodology is strict context
isolation, under which prompt refinement and content generation are treated
as entirely separate operations, each carried out without knowledge of the
other's context.

\section{Architect Chat}
\label{AIPromptEngineering:ArchitectChat}

An initial ``Architect Chat'', run with Gemini Pro 3.1 Extended, was used to
plan and outline the optimal sequence of prompts, ensuring that all project
requirements were systematically addressed before any content generation
began. This planning stage produced the list of prompts later refined and
executed in the two phases described below, but did not itself generate any
content included in the report.

\section{Prompt Engineering Phases}
\label{AIPromptEngineering:Phases}

\subsection{Prompt Improvement}
\label{AIPromptEngineering:Phases:PromptImprovement}

The first phase focuses on prompt refinement. During this stage, the
initial prompts produced by the Architect Chat are iteratively rewritten to
align with the project guidelines and to maintain a rigorous academic
standard. Each prompt is refined independently: the AI model is tasked
solely with improving the prompt's structure and clarity, without executing
it and without carrying over context from unrelated topics or previous
prompts.

\subsection{Prompt Execution}
\label{AIPromptEngineering:Phases:PromptExecution}

The second phase involves executing the refined prompts to generate the
final content, using Claude Sonnet 5 in a new, isolated session for each
prompt. This isolation ensures that the generated content relies
exclusively on the explicit instructions contained in the finalized prompt,
without any context leakage from the improvement phase or from other
prompts. During execution, the model is also explicitly instructed to
report the source references and documentation used to formulate its
response.

\section{Prompt 0: Configure Prompt Engineering (Task 1)}
\label{Prompt 0 Task1}

To ensure the generation of rigorous academic sections, this methodology
treats prompt generation and the final report generation as strictly independent
operations. The process does not rely on manual, ad-hoc prompts for each
language. Instead, it employs a \textit{meta-prompting} approach. A
\textbf{Master Prompt} was designed whose sole purpose is to act as a
"Senior Prompt Engineer" in an isolated session. This system takes a raw
specification as input and outputs a formal, optimized execution prompt,
which is then used in a completely separate session to generate the final text.

\noindent\textbf{Master Prompt Specification} \\
Below is the complete text of the Master Prompt used to orchestrate
the prompt engineering phase for Task 1 (language implementation of parametric
polymorphism):

\begin{PromptVerbatim}
System Persona:

Act as a Senior Prompt Engineer specialized in designing rigorous prompts
for the generation of technical-academic reports in computer science,
programming language theory, compiler design, and software engineering.

Goal:

Given, as input, a task specification containing a persona, goal,
requirements, writing-style constraints, output structure, and possible
reference requirements for producing a technical-academic report, generate a
single, self-contained master prompt ready to be executed in a new language
model conversation.

The generated master prompt must preserve the original intent and must be
optimized to produce a formally structured academic report with appropriate
technical depth, logical organization, and scientific writing standards.

Requirements:

Task 1, Extraction and Preservation:

Parse the input specification and identify:
- The requested system persona or expert role.
- The report objective and title, if provided.
- All explicit and implicit technical requirements.
- Required theoretical concepts, comparisons, analyses, and evaluations.
- Writing-style constraints.
- Required output format and document structure.
- Requirements related to references, citations, bibliography, or external
  academic sources.

Preserve the complete meaning of the original specification.
Do not remove requirements.
Do not introduce new technical requirements.
Do not change the intended scope of the report.
Do not simplify technical constraints.

Task 2, Correction and Structuring:

Correct grammatical errors, ambiguities, inconsistent terminology,
formatting problems, and unclear instructions present in the input.

Reorganize the extracted information into a clear and logically structured
master prompt using the following sections:
- System Persona
- Goal
- Requirements
- Writing Style Constraint
- Output Structure

When reference requirements are present, preserve them explicitly in the
Requirements section.

Task 3, Academic Report Generation Constraints:

Ensure that the generated master prompt instructs the target language model
to:
- Produce a complete technical-academic report rather than a short summary.
- Maintain a formal, objective, and scientific writing style.
- Use precise terminology appropriate for computer science literature.
- Provide a logical progression from theoretical foundations to technical
  analysis.
- Include comparisons, trade-offs, and implementation details when required.
- Distinguish clearly between theoretical concepts and implementation
  mechanisms.
- Avoid unsupported claims and clearly identify assumptions.
- ENFORCE STRICT SCOPE CONTROL: Explicitly instruct the target language model
  to confine the generated text strictly to the requested core topic. The model
  must actively avoid tangential subjects (e.g., memory footprint, hardware
  performance, code bloat, garbage collection overhead, or broader language
  comparisons) unless they are explicitly demanded by the input specification.

Task 4, Reference and Citation Requirements:

If the original specification requires the use of references:
- Require the generated report to rely on authoritative academic or technical
  sources.
- Prefer peer-reviewed publications, academic books, official language
  specifications, and authoritative technical documentation.
- Require consistent citation formatting.
- Require a final bibliography section containing all cited resources.
- Avoid fabricated references, unverifiable sources, or unsupported citations.

Task 5, LaTeX Output Requirements:

If the original specification requires LaTeX output:
- Require complete compilation-ready LaTeX code.
- Preserve a professional academic document structure.
- Include appropriate sections, chapters, mathematical notation, tables,
  figures, citations, and bibliography commands when required.
- Ensure compatibility with standard LaTeX compilation workflows.

Writing Style Constraint (STRICT):

Do not execute the input specification.
Do not generate the final report.
Do not answer the original technical question.

The output must contain only the generated master prompt.

Do not include:
- explanations,
- comments,
- editorial notes,
- placeholders,
- analysis of the transformation,
- meta-commentary addressed to the user.

Do not use first-person forms ("I", "we") in the generated master prompt unless
they are explicitly required by the original specification.

Output Structure:

Return only the generated master prompt.

The generated prompt must be formatted as ready-to-use text and organized
using the following headings:

System Persona

Goal

Requirements

Writing Style Constraint

Output Structure

No additional text must appear before or after the generated master prompt.
\end{PromptVerbatim}

\section{Prompt 0bis: Configure Prompt Engineering (Task 2)}
\label{Prompt 0 Task2}

For Task 2 (metaprogramming across C++, Java, and C\#), a second Master
Prompt was used. Its purpose is identical to the one described above, but it
is specialized to produce a single \emph{section} of an existing chapter
(rather than a standalone report), given one focused instruction as input.

\begin{PromptVerbatim}
System Persona:

Act as a Senior Prompt Engineer specialized in designing rigorous prompts
for the generation of individual sections within a larger technical-
academic report chapter on programming language theory and software
engineering.

Goal:

Given, as input, a single focused instruction (below, marked as
[INPUT PROMPT]), generate a single, self-contained master prompt ready to
be executed in a new language model conversation. The master prompt must
instruct the target language model to produce ONE SECTION of a larger
chapter -- not a standalone report -- on parametric polymorphism and
metaprogramming across C++, Java, and C#.

[INPUT PROMPT]:
"{{PASTE HERE ONE OF THE 9 PROMPTS}}"

Requirements:

Task 1 -- Extraction and Preservation:

Parse [INPUT PROMPT] and identify:
- The specific technical question or comparison it asks for.
- The languages, techniques, or mechanisms explicitly named.
- Any implicit expectation (e.g. code examples, explanation of underlying
  mechanism, comparison table).

Preserve the complete meaning of [INPUT PROMPT]. Do not remove
requirements. Do not introduce new technical requirements beyond what
[INPUT PROMPT] asks. Do not expand the scope beyond the specific question
posed.

Task 2 -- Section-Level Framing (STRICT):

The generated master prompt must explicitly instruct the target language
model that its output is ONE SECTION of an existing chapter, not a
complete report. Specifically, the master prompt must require the target
model to:
- Omit any standalone introduction, motivation, or "in this report we will"
  framing -- begin directly with the technical content.
- Omit any standalone conclusion, summary of the whole chapter, or
  forward-looking remarks about other sections.
- Omit its own bibliography or reference list -- citations, if any, are
  consolidated once at the end of the full chapter, not per section.
- Assume the reader already has the context of the broader comparison
  (C++ templates vs Java generics vs C# generics) and does not need it
  re-introduced.

Task 3 -- Content Constraints:

Ensure the generated master prompt instructs the target language model to:
- Maintain a formal, objective, and precise technical writing style
  consistent with programming language theory literature.
- Use precise terminology (instantiation, monomorphization, type erasure,
  reification, substitution failure, bounded polymorphism) only where
  relevant to [INPUT PROMPT].
- Include working code examples only if [INPUT PROMPT] explicitly or
  implicitly requires them.
- Distinguish clearly between compile-time and runtime mechanisms when
  relevant to the specific question.
- Avoid unsupported claims; state assumptions about language/compiler
  version only if relevant to the specific question asked.
- ENFORCE STRICT SCOPE CONTROL: address only what [INPUT PROMPT] asks.
  Do not digress into memory footprint, hardware performance benchmarking,
  code bloat, garbage collection overhead, or unrelated language
  comparisons unless [INPUT PROMPT] explicitly requires it.
- ENFORCE CONCISENESS: the section must be as short as fully answering
  [INPUT PROMPT] allows. No padding, no restating the question, no
  redundant transitions between paragraphs.

Writing Style Constraint (STRICT):

Do not execute [INPUT PROMPT] directly. Do not generate the final section
content. Do not answer the original technical question. The output must
contain only the generated master prompt. Do not include explanations,
comments, editorial notes, placeholders, analysis of the transformation,
or meta-commentary addressed to the user. Do not use first-person forms
("I", "we") in the generated master prompt unless explicitly required by
[INPUT PROMPT].

Output Structure:

Return only the generated master prompt, formatted as ready-to-use text
and organized using the following headings: System Persona, Goal,
Requirements, Writing Style Constraint, Output Structure. No additional
text must appear before or after the generated master prompt.
\end{PromptVerbatim}

% =====================================================================
\chapter{Task 1 -- Language Implementation of Parametric Polymorphism}
% =====================================================================

This chapter groups the prompts used to generate the report content
addressing Task 1: \textit{``Explain how each language implements parametric
polymorphism (C++ templates, Java generics, C\# generics), including when
instantiation happens, what exists at runtime, and how this affects
performance, type safety, and reflection.''} Each section below merges what
were previously three separate, per-language prompts (C++, Java, C\#) into a
single comparative prompt covering all three languages together, so that a
single generation directly produces the cross-language comparison.

\vspace{0.5cm}
\noindent\textbf{Methodological Note on Prompt Construction and Verification:}\\
To ensure the accuracy of the generated analysis, the prompt engineering strategy utilized a decoupled approach. Specifically,
individual prompts were tailored for each language (C++, Java, and C\#) across five core analytical categories:
\begin{enumerate}
    \item Timing of instantiation
    \item Runtime representation
    \item Performance implications
    \item Type safety
    \item Reflection capabilities
\end{enumerate}

Because multiple distinct queries targeted different facets of the same language, the model's outputs naturally contained
redundant technical information. This redundancy was intentional; it served as the critical first step in the verification
process, allowing for the cross-referencing of the model's overlapping responses to ensure internal consistency before
synthesizing the final master prompt.

To simplify the analysis and improve readability within this document, these isolated queries were subsequently merged.
The resulting unified master prompt dictates the exact structure,
constraints, and content parameters used to
generate the comprehensive comparison presented in the following sections.

\noindent\textit{Note on repeated boilerplate:} from Prompt 3 onward, the
\textbf{Writing Style Constraint} and \textbf{Output Structure} blocks follow
the same template established in Prompt 2 (\S\ref{Prompt 2 Instantiation}),
adapted only in terminology and section titles to the specific topic. To avoid
redundancy, these two blocks are stated in full only once (Prompt 2) and
referenced thereafter.

% =====================================================================
\section{Prompt 1: Overview of Parametric Polymorphism}
% =====================================================================

\noindent\textbf{Model:} ChatGPT-5.5

\vspace{0.5em}
\noindent\textbf{Goal:} Generate an overview of Parametric Polymorphism and ad-hoc polymorphism, giving details of type erasure compared with generics and monomorphization, and comparing differences between source code, bytecode, and binary code.

\vspace{0.5em}
\noindent\textbf{Prompt}:

\begin{PromptVerbatim}
System Persona:

Act as a world-class computer scientist and academic researcher specialized
in programming language theory, type systems, and the historical foundations
of polymorphism.

Goal:

Produce an exhaustive, publication-grade theoretical overview of parametric
polymorphism and ad-hoc polymorphism, comparing type erasure against
reified generics and monomorphization, and explaining how these strategies
manifest differently across source code, bytecode, and binary code.

Requirements:

- Provide an exhaustive theoretical background on parametric polymorphism.
- Define and thoroughly explain Christopher Strachey's distinction between
  parametric polymorphism and ad-hoc polymorphism.
- Detail the two primary implementation strategies discussed in the
  literature:
  1. Homogeneous translation (type erasure), detailing how it results in
     one compiled representation for all instantiations.
  2. Heterogeneous translation (reification/monomorphization), detailing
     how it results in separate code generated per instantiated type.
- Analyze type erasure versus reified generics as the core axis of
  comparison.
- Analyze compile-time versus runtime characteristics as the second axis
  of comparison, specifically examining when instantiation happens and
  what type information survives into the compiled binary or bytecode.

Writing Style Constraint:

Maintain a formal, objective, and precise academic writing style consistent
with programming language theory literature. Avoid first-person pronouns.
Provide technically detailed explanations rather than high-level summaries.

Output Structure:

The entire document must be written in professional, compilation-ready
LaTeX code, organized into clearly labeled sections covering the
theoretical background, the Strachey distinction, the two translation
strategies, and the comparative analysis.
\end{PromptVerbatim}

\vspace{0.5em}
\noindent\textbf{Sources}:
\begin{itemize}
    \item \textbf{Strachey, ``Fundamental Concepts in Programming Languages''} \cite{strachey1967} -- used for the original classification of parametric and ad-hoc polymorphism.
    \item \textbf{Cardelli and Wegner, ``On Understanding Types, Data Abstraction, and Polymorphism''} \cite{cardelli1985} -- used for the analysis of type abstraction and forms of polymorphism.
    \item \textbf{Pierce, \emph{Types and Programming Languages}} \cite{pierce2002} -- used for the formal treatment of parametric polymorphism and type systems.
    \item \textbf{Pierce, \emph{Advanced Topics in Types and Programming Languages}} \cite{pierce2005} -- used for advanced insights into type systems.
    \item \textbf{Kennedy and Syme, ``Design and Implementation of Generics for the .NET CLR''} \cite{kennedy2001} -- used for the description of runtime reification in the CLR.
    \item \textbf{Bracha et al., ``Making the Future Safe for the Past''} \cite{bracha1998} -- used for the type erasure mechanism introduced in Java generics.
    \item \textbf{Bracha, \emph{Generics in the Java Programming Language}} \cite{bracha2004} -- used for the technical documentation of Java erasure.
    \item \textbf{ECMA-334: C\# Language Specification} \cite{ecma334_2023} -- used as a normative reference for C\# generics.
    \item \textbf{ECMA-335: Common Language Infrastructure} \cite{ecma335_2024} -- used for the runtime representation of generics in the CLR.
    \item \textbf{The Java Language Specification} \cite{jls2024} -- used as a normative reference for Java generics.
    \item \textbf{The Java Virtual Machine Specification} \cite{jvms2024} -- used for the bytecode representation resulting from erasure.
    \item \textbf{Cooper and Torczon, \emph{Engineering a Compiler}} \cite{cooper2022} -- used for general compilation context.
    \item \textbf{Appel, \emph{Modern Compiler Implementation}} \cite{appel1998} -- used for general compilation context.
    \item \textbf{Aho et al., \emph{Compilers: Principles, Techniques, and Tools}} \cite{aho2022} -- used for compiler theory.
    \item \textbf{Official Rust documentation} \cite{rust2024} -- used for a comparison on monomorphization in other languages.
    \item \textbf{Rust documentation on monomorphization} \cite{rust_monomorphization} -- used to detail the monomorphization mechanism.
    \item \textbf{Vandevoorde, Josuttis, Gregor, \emph{C++ Templates: The Complete Guide}} \cite{cpp_templates2017} -- used for the C++ template instantiation model.
    \item \textbf{Swift documentation on generics} \cite{swift_generics2019} -- used for an additional comparison on generics implementation strategies.
\end{itemize}

% =====================================================================
\section{Prompt 2: Instantiation Time -- C++, Java, and C\# Compared}
\label{Prompt 2 Instantiation}
% =====================================================================

\noindent\textbf{Goal:} Generate a single, unified, technically rigorous section comparing the \emph{instantiation time} of parametric polymorphism across C++ templates, Java generics, and C\# generics, covering compile-time template instantiation, compile-time type erasure, and runtime CLR/JIT-based instantiation within one comparative treatment.

\noindent\textbf{Model:} ChatGPT-5.5

\vspace{0.5em}

\vspace{0.5em}
\noindent\textbf{Prompt}:

\begin{PromptVerbatim}
System Persona:

Act as a world-class computer scientist and principal compiler/runtime
engineer specializing in Programming Language Theory (PLT), type systems,
and the compiler/virtual-machine execution models of C++, the JVM, and the
.NET CLR.

Goal:

Generate an exhaustive, publication-grade academic report section
analyzing the instantiation time of parametric polymorphism, contrasting
three distinct temporal models: C++ compile-time template instantiation,
Java compile-time type erasure, and C# runtime CLR/JIT instantiation.

Requirements:

- C++ instantiation model: Explain the complete compile-time instantiation
  pipeline, including template argument deduction, two-phase name lookup,
  implicit and explicit template instantiation, and monomorphization/
  specialization. Emphasize that no template instantiation occurs during
  program execution. Describe how multiple translation units may
  independently instantiate identical specializations and how linkers
  eliminate duplicates (COMDAT sections, weak symbols).
- Java instantiation model: Explain the Java Generics compilation
  pipeline, including compile-time verification of generic constraints
  and bounds, type checking of parameterized types, type erasure
  transformation, replacement of type variables with erasures, generation
  of synthetic bridge methods, and insertion of implicit runtime casts in
  the generated bytecode. Emphasize that no generic instantiation occurs
  at runtime; only a single erased representation exists.
- C# instantiation model: Explain the dual-phase mechanism of C# generics:
  compile-time emission of Common Intermediate Language (CIL) with generic
  metadata preserved, followed by runtime Just-In-Time (JIT) instantiation.
  Detail reference-type code sharing versus value-type specialization
  performed by the JIT.
- Comparative synthesis: Directly compare the three timelines -- C++
  ahead-of-time compile-time monomorphization, Java compile-time erasure
  with no runtime instantiation, and C# runtime JIT-time specialization --
  making explicit when instantiation happens in each model and what
  triggers it.
- Whenever appropriate, relate each implementation strategy to the
  theoretical notion of type substitution in parametric polymorphism.
- ENFORCE STRICT SCOPE CONTROL: confine the text strictly to instantiation
  timing and mechanisms. Do not discuss executable size, code bloat,
  memory footprint, garbage collection overhead, or general runtime
  performance/speed.
- Base every technical statement on authoritative references, prioritizing
  the ISO/IEC 14882 C++ Standard, the Java Language Specification, the
  Java Virtual Machine Specification, ECMA-334, and ECMA-335.
- Include a complete bibliography using standard LaTeX thebibliography
  commands.

Writing Style Constraint:

- Adopt a formal, objective, and precise scientific writing style
  comparable to ACM or IEEE publications.
- Use accurate compiler/runtime terminology (e.g., template argument
  deduction, two-phase lookup, monomorphization, translation unit,
  COMDAT, type erasure, bridge method, JIT specialization, reification).
- Avoid first-person pronouns and conversational language.
- Provide technically detailed explanations rather than high-level
  summaries.

Output Structure:

The generated document must consist entirely of valid, self-contained,
compilation-ready LaTeX code, organized as:

1. Introduction: Three Models of Instantiation Time
2. C++ Templates: Compile-Time Instantiation Pipeline
3. Java Generics: Compile-Time Type Erasure
4. C# Generics: Runtime CLR/JIT Instantiation
5. Comparative Timeline: When Instantiation Happens in Each Language
6. Conclusion
\end{PromptVerbatim}

\vspace{0.5em}
\noindent\textbf{Sources}:
\begin{itemize}
    \item \textbf{ISO/IEC 14882:2020/2024, Programming Languages --- C++} \cite{iso_cpp2024} -- used as a normative reference for template instantiation.
    \item \textbf{Vandevoorde, Josuttis, Gregor, \emph{C++ Templates: The Complete Guide}} \cite{cpp_templates2017} -- used for the template deduction and instantiation pipeline.
    \item \textbf{Stroustrup, \emph{The C++ Programming Language}} \cite{stroustrup2013} -- used for the general C++ compilation model.
    \item \textbf{The Java Language Specification, Java SE 24} \cite{jls2024} -- used for the type erasure mechanism.
    \item \textbf{The Java Virtual Machine Specification, Java SE 24} \cite{jvms2024} -- used for the bytecode representation resulting from erasure.
    \item \textbf{Bracha et al., ``Making the Future Safe for the Past''} \cite{bracha1998} -- used for the introduction of generics in Java.
    \item \textbf{Bracha, \emph{Generics in the Java Programming Language}} \cite{bracha2004} -- used for the implementation details of erasure.
    \item \textbf{Igarashi, Pierce, Wadler, ``Featherweight Java''} \cite{igarashi2001} -- used for the formal model of Java generics.
    \item \textbf{ECMA-334: C\# Language Specification} \cite{ecma334_2023} -- used as a normative reference for C\# generics.
    \item \textbf{ECMA-335: Common Language Infrastructure} \cite{ecma335_2024} -- used for the JIT instantiation mechanism.
    \item \textbf{Kennedy and Syme, ``Design and Implementation of Generics for the .NET CLR''} \cite{kennedy2001} -- used for the JIT sharing/specialization mechanism.
    \item \textbf{Richter, \emph{CLR via C\#}} \cite{richter2012} -- used for the runtime behavior of the CLR.
    \item \textbf{Pierce, \emph{Types and Programming Languages}} \cite{pierce2002} -- used for the theoretical concept of type substitution.
    \item \textbf{Aho et al., \emph{Compilers: Principles, Techniques, and Tools}} \cite{aho2022} -- used for general compilation context.
    \item \textbf{Cooper and Torczon, \emph{Engineering a Compiler}} \cite{cooper2022} -- used for general compilation context.
\end{itemize}

% =====================================================================
\section{Prompt 3: Runtime Representation -- C++, Java, and C\# Compared}
\label{Prompt 3 Runtime}
% =====================================================================

\noindent\textbf{Goal:} Generate a single, unified section analyzing the \emph{runtime representation} of parametric polymorphism across C++ templates (monomorphized native code, no generic metadata), Java generics (erased representation), and C\# generics (reified CLR representation).

\noindent\textbf{Model:} ChatGPT-5.5

\vspace{0.5em}
\noindent\textbf{Prompt}:

\begin{PromptVerbatim}
System Persona:

Act as a world-class computer scientist and principal compiler/runtime
engineer specializing in Programming Language Theory (PLT), compiler
backends (GCC/Clang), the Java Virtual Machine, and the .NET Common
Language Runtime.

Goal:

Generate an exhaustive, publication-grade academic report section
analyzing the instantiation time of parametric polymorphism, contrasting
three distinct temporal models: C++ compile-time template instantiation,
Java compile-time type erasure, and C# runtime CLR/JIT instantiation.

Requirements:

- C++: Explain heterogeneous translation (monomorphization) as producing
  distinct, independent native machine code artifacts per instantiated
  type. Clarify the absence of generic metadata at runtime: the execution
  environment sees only ordinary functions and objects, not templates.
  Address the limited type information that survives (name mangling,
  RTTI) and how this differs from reification.
- Java: Explain how type erasure produces a single shared runtime
  representation with no generic metadata for parameter types, the
  insertion of implicit casts, bridge methods, and the consequences for
  primitive types (boxing/unboxing).
- C#: Explain CLR reification: how generic metadata is preserved in CIL
  and how the JIT constructs runtime representations, distinguishing
  reference-type code sharing (with runtime dictionary passing) from
  value-type specialization (distinct native code per value type).
- Comparative synthesis: Directly contrast the three runtime
  representations along the axis of metadata presence and code-artifact
  granularity: C++ (no metadata, per-type native code), Java (shared
  erased representation, no metadata), C# (preserved metadata, partial
  code sharing).
- ENFORCE STRICT SCOPE CONTROL: confine the text strictly to runtime
  representation and memory/metadata structure. Do not discuss
  compilation times, code-bloat metrics, manual memory management, or
  CPU cache performance.
- Base every technical statement on authoritative references (ISO C++
  Standard, JLS, JVMS, ECMA-335, and established compiler/runtime
  literature).
- Include a complete bibliography using standard LaTeX thebibliography
  commands.

Writing Style Constraint:

As defined for the instantiation-time section, with terminology adapted
to runtime representation (monomorphization, name mangling, RTTI, type
erasure, bridge method, reification, generic dictionary, CIL metadata).

Output Structure:

1. Introduction: Three Models of Runtime Representation
2. C++: Monomorphized Native Code and Absence of Generic Metadata
3. Java: Erased Shared Representation
4. C#: Reified CLR Representation and JIT-Constructed Types
5. Comparative Analysis: Metadata Presence and Code-Artifact Granularity
6. References
\end{PromptVerbatim}

\vspace{0.5em}
\noindent\textbf{Sources}:
\begin{itemize}
    \item \textbf{Vandevoorde, Josuttis, Gregor, \emph{C++ Templates: The Complete Guide}} \cite{cpp_templates2017} -- used for the monomorphization model.
    \item \textbf{ISO/IEC 14882:2024, Programming Languages --- C++} \cite{iso_cpp2024} -- used as a normative reference.
    \item \textbf{Stroustrup, \emph{The C++ Programming Language}} \cite{stroustrup2013} -- used for the C++ object model.
    \item \textbf{Lippman, \emph{Inside the C++ Object Model}} \cite{lippman1996} -- used for the internal representation of C++ objects.
    \item \textbf{Itanium C++ ABI Specification} \cite{itanium_abi} -- used for name mangling and RTTI.
    \item \textbf{The Java Language Specification, Java SE 24} \cite{jls2024} -- used for the erasure mechanism.
    \item \textbf{The Java Virtual Machine Specification, Java SE 24} \cite{jvms2024} -- used for the bytecode representation.
    \item \textbf{Bracha, \emph{Generics in the Java Programming Language}} \cite{bracha2004} -- used for bridge methods and boxing.
    \item \textbf{ECMA-335: Common Language Infrastructure} \cite{ecma335_2024} -- used for CLR reification.
    \item \textbf{Kennedy and Syme, ``Design and Implementation of Generics for the .NET CLR''} \cite{kennedy2001} -- used for reference-type/value-type code sharing.
    \item \textbf{Richter, \emph{CLR via C\#}} \cite{richter2012} -- used for JIT runtime behavior.
\end{itemize}

% =====================================================================
\section{Prompt 4: Performance Characteristics -- C++, Java, and C\# Compared}
\label{Prompt 4 Performance}
% =====================================================================

\noindent\textbf{Goal:} Generate a single, unified section analyzing the runtime performance implications of each language's implementation strategy, following directly from the runtime representation discussion.

\noindent\textbf{Model:} ChatGPT-5.5

\vspace{0.5em}
\noindent\textbf{Prompt}:

\begin{PromptVerbatim}
System Persona:

Act as a Senior Computer Science Researcher and Expert in Programming
Language Theory, Compiler Design, JVM Architecture, and .NET CLR
Architecture.

Goal:

Generate a comprehensive, technically rigorous academic section analyzing
and directly comparing the execution performance characteristics of C++
templates (compile-time monomorphization), Java generics (type erasure),
and C# generics (runtime reification with JIT specialization).

Requirements:

- Establish continuity with the following preceding theoretical context:
  "Both implementation strategies preserve static type safety, although
  they achieve this objective differently. In implementations based on
  type erasure, generic correctness is verified entirely during
  compilation. Although type parameters disappear from the generated
  executable, compiler-inserted casts preserve the guarantees established
  by the static type system. Implementations based on reified generics
  retain complete type information throughout compilation and generate
  type-specific implementations directly."
- C++: Detail the performance advantages of monomorphization: static
  dispatch, elimination of boxing/unboxing, and facilitation of
  aggressive compiler optimizations such as inlining.
- Java: Detail the runtime overheads introduced by type erasure: pointer
  indirection, mandatory boxing/unboxing for primitive types, impacts on
  heap allocation and data locality. Discuss JIT-level mitigations
  (escape analysis, inline caching, devirtualization) only insofar as
  they mitigate generic-specific overhead.
- C#: Detail the distinct performance behavior of value-type
  specialization (per-type native code, no boxing) versus reference-type
  sharing (single implementation, dictionary-passed type handles) under
  JIT compilation.
- Comparative synthesis: Directly compare static dispatch/AOT
  optimization (C++), erasure-induced boxing overhead (Java), and the
  hybrid value/reference-type performance profile (C#).
- ENFORCE STRICT SCOPE CONTROL: confine the analysis to execution
  performance, runtime throughput, memory layout, and cache behavior. Do
  not digress into compilation times, language ergonomics, or general
  garbage collection mechanics unless directly tied to generic-structure
  throughput.
- Rely on authoritative sources: ISO C++ Standard, JLS, JVMS, ECMA-335,
  C# Language Specification, and established compiler literature.
  Include consistent citations and a final bibliography.

Writing Style Constraint:

As defined for the instantiation-time section, with terminology adapted
to performance analysis (static dispatch, boxing overhead, data locality,
reification, monomorphization, type erasure, JIT specialization, cache
lines).

Output Structure:

1. Introduction: Performance as a Consequence of Implementation Strategy
2. C++: Static Dispatch and Compile-Time Optimization
3. Java: Type Erasure and Boxing Overhead
4. C#: Value-Type Specialization vs. Reference-Type Sharing
5. Comparative Analysis: Throughput, Memory Layout, and Cache Behavior
6. References
\end{PromptVerbatim}

\vspace{0.5em}
\noindent\textbf{Sources}:
\begin{itemize}
    \item \textbf{Vandevoorde, Josuttis, Gregor, \emph{C++ Templates: The Complete Guide}} \cite{cpp_templates2017} -- used for the performance advantages of monomorphization.
    \item \textbf{Muchnick, \emph{Advanced Compiler Design and Implementation}} \cite{muchnick1997compiler} -- used for compiler optimizations (inlining, static dispatch).
    \item \textbf{The Java Language Specification, Java SE 21/24} \cite{jls2024} -- used for boxing/unboxing behavior.
    \item \textbf{ECMA-335: Common Language Infrastructure} \cite{ecma335_2024} -- used for JIT specialization of value types.
    \item \textbf{Microsoft Corporation, C\# Language Specification} \cite{CSharpSpec} -- used as a normative reference for C\# generics.
\end{itemize}

% =====================================================================
\section{Prompt 5: Type Safety -- C++, Java, and C\# Compared}
\label{Prompt 5 TypeSafety}
% =====================================================================

\noindent\textbf{Goal:} Generate a single, unified section analyzing static type safety in C++ templates, Java generics, and C\# generics, illustrated with explicative code examples for each language.

\noindent\textbf{Model:} ChatGPT-5.5

\vspace{0.5em}
\noindent\textbf{Prompt}:

\begin{PromptVerbatim}
System Persona:

Act as an expert academic researcher and computer scientist specializing
in programming language theory, compiler design, and cross-language
generic type-system semantics.

Goal:

Generate an exhaustive, academic-grade LaTeX section analyzing static
type safety across C++ templates (reified generics via monomorphization),
Java generics (type erasure), and C# generics (reified generics via CLR).

Requirements:

- Integrate the following theoretical premise as the foundation: "Both
  implementation strategies preserve static type safety, although they
  achieve this objective differently. In implementations based on type
  erasure, generic correctness is verified entirely during compilation.
  Although type parameters disappear from the generated executable,
  compiler-inserted casts preserve the guarantees established by the
  static type system. Implementations based on reified generics retain
  complete type information throughout compilation and generate
  type-specific implementations directly. Since each specialization
  corresponds to a concrete type, correctness is preserved without
  requiring erased representations or implicit runtime casts."
- C++: Detail how templates enforce type safety at compile time via
  monomorphization, retaining complete type information, with a simple
  explanatory code example.
- Java: Detail how type erasure verifies generic correctness entirely at
  compile time via compiler-inserted casts, with a simple explanatory
  code example.
- C#: Detail how reified generics retain complete type information
  throughout compilation and generate type-specific implementations, with
  a simple explanatory code example.
- Relate each code example directly to the theoretical concepts of
  reification, erasure, and concrete type specialization.
- Rely on authoritative academic or technical sources (ISO/IEC 14882,
  JLS, ECMA-335, and peer-reviewed PLT literature).
- ENFORCE STRICT SCOPE CONTROL: confine the text strictly to type-safety
  mechanisms and the theoretical strategies of implementation. Do not
  digress into memory footprint, hardware performance, code bloat,
  compilation times, or garbage collection overhead.

Writing Style Constraint:

As defined for the instantiation-time section, adapted to type safety,
distinguishing theoretical concepts (parametric polymorphism) from
implementation mechanisms (monomorphization, erasure, reification), with
code examples formatted using the standard listings package.

Output Structure:

1. Introduction: Type Safety Under Different Implementation Strategies
2. C++: Type Safety via Reified Templates
3. Java: Type Safety via Type Erasure
4. C#: Type Safety via CLR Reification
5. Comparative Synthesis
6. References
\end{PromptVerbatim}

\vspace{0.5em}
\noindent\textbf{Sources}:
\begin{itemize}
    \item \textbf{Vandevoorde, Josuttis, Gregor, \emph{C++ Templates: The Complete Guide}} \cite{cpp_templates2017} -- used for the type-safety model of templates.
    \item \textbf{Stroustrup, \emph{The C++ Programming Language}} \cite{stroustrup2013} -- used for static type checking in C++.
    \item \textbf{Pierce, \emph{Types and Programming Languages}} \cite{pierce2002} -- used for the theoretical treatment of type safety.
    \item \textbf{Gosling et al., The Java Language Specification} \cite{gosling2014java} -- used for the mechanism of compiler-inserted casts.
    \item \textbf{Microsoft Corporation, C\# Language Specification} \cite{CSharpSpec} -- used for the reification model of C\# generics.
    \item \textbf{ECMA-335: Common Language Infrastructure} \cite{ecma335_2024} -- used for runtime type verification in the CLR.
\end{itemize}

% =====================================================================
\section{Prompt 6: Reflection and Runtime Type Information -- C++, Java, and C\# Compared}
\label{Prompt 6 Reflection}
% =====================================================================

\noindent\textbf{Goal:} Generate a single, unified section analyzing reflection and runtime type information (RTTI) across C++ templates, Java generics, and C\# generics.

\noindent\textbf{Model:} ChatGPT-5.5

\vspace{0.5em}
\noindent\textbf{Prompt}:

\begin{PromptVerbatim}
System Persona:

Act as a Senior Research Scientist in Programming Language Theory,
Compiler Architecture, and Type Systems, with expertise in C++ RTTI, Java
reflection, and .NET System.Reflection.

Goal:

Generate a rigorous, technical-academic report section analyzing
reflection and RTTI capabilities across C++, Java, and C#, grounded in the
theoretical premise that runtime availability of generic type information
depends directly on the chosen representation strategy: type-erased
implementations discard most generic information, while reified
implementations preserve it for runtime introspection.

Requirements:

- C++: Detail the application and behavior of standard RTTI mechanisms
  (typeid, std::type_info, dynamic_cast) on instantiated template
  classes/functions, and how monomorphization preserves concrete type
  information per specialization.
- Java: Detail how type erasure constrains reflection: why instanceof
  against parameterized types is disallowed, why getClass() yields an
  identical reference across structurally distinct generic
  instantiations, and how the type-token pattern (Class<T>) is used as a
  workaround. Include a compilation-ready Java code example.
- C#: Detail how CLR reification preserves generic metadata for full
  runtime introspection via System.Reflection, including IsGenericType
  and GetGenericArguments(). Include a compilation-ready C# code example.
- Comparative synthesis: Directly contrast the three reflection models
  along the axis of metadata availability: C++ (per-specialization RTTI,
  no generic-level metadata), Java (single shared Class object, no
  parameter recovery without workarounds), C# (full runtime introspection
  of type arguments).
- ENFORCE STRICT SCOPE CONTROL: confine the text strictly to reflection,
  metadata availability, and type introspection capabilities. Do not
  discuss memory footprint, hardware performance, code bloat, or garbage
  collection overhead.
- Base every technical statement on authoritative references and avoid
  fabricated citations.

Writing Style Constraint:

As defined for the instantiation-time section, adapted to reflection and
RTTI terminology (RTTI, type erasure, reification, type token,
monomorphization, generic metadata).

Output Structure:

1. Introduction: Reflection as a Function of Representation Strategy
2. C++: RTTI on Monomorphized Templates
3. Java: Type Erasure and Reflection Limitations
4. C#: Full Runtime Introspection via CLR Reification
5. Comparative Analysis
6. References
\end{PromptVerbatim}

\vspace{0.5em}
\noindent\textbf{Sources}:
\begin{itemize}
    \item \textbf{Stroustrup, \emph{The C++ Programming Language}} \cite{stroustrup2013} -- used for the RTTI mechanism in C++.
    \item \textbf{Vandevoorde, Josuttis, Gregor, \emph{C++ Templates: The Complete Guide}} \cite{cpp_templates2017} -- used for RTTI applied to instantiated templates.
    \item \textbf{cppreference.com, C++ Language Reference: RTTI and typeid} \cite{cppreference} -- used for the syntax and behavior of typeid and dynamic\_cast.
    \item \textbf{Oracle, The Java Language Specification} \cite{jls2024} -- used for the limitations of instanceof on parameterized types.
    \item \textbf{Oracle, The Java Virtual Machine Specification} \cite{jvms2024} -- used for the behavior of getClass() under erasure.
    \item \textbf{Richter, \emph{CLR via C\#}} \cite{richter2012} -- used for introspection via System.Reflection in the CLR.
\end{itemize}

% =====================================================================
\chapter{Task 2 -- Metaprogramming Across C++, Java, and C\#}
% =====================================================================

This chapter groups the prompts used to generate the report content
addressing Task 2: \textit{``Provide examples of metaprogramming in C++
using templates (e.g. compile-time computation, SFINAE / type traits) and
compare them with what is and isn't possible using Java and C\# generics,
illustrating the practical limits of metaprogramming in each ecosystem.''}
The sections retain their existing individual structure -- C++-native
techniques are treated first, followed by dedicated sections analyzing the
limits of Java and of C\#, and finally the workarounds used in Java and C\#
to approximate C++ metaprogramming.

\noindent\textit{Note on repeated boilerplate:} all prompts in this chapter
share the same \textbf{Section-Level Framing} constraints (output is a single
section of an existing chapter; no standalone introduction, conclusion, or
bibliography; formal PLT writing style; strict scope control; enforced
conciseness), stated in full in Prompt 1 and referenced thereafter rather than
repeated verbatim.

% =====================================================================
\section{Prompt 1: Concrete C++ Template Examples of Compile-Time Computation}
% =====================================================================

\noindent\textbf{Goal}: Generate a detailed academic section of compile-time computation in C++ templates. For each example, explain what the compiler actually generates and why the computation happens entirely before runtime. Include full compilable code.

\noindent\textbf{Model:} ChatGPT-5.5

\noindent\textbf{Prompt}:
\begin{PromptVerbatim}
System Persona:

Act as an expert in programming language theory and advanced C++ software engineering.

Goal:

Write a single, highly focused section of a larger technical-academic chapter discussing parametric polymorphism and metaprogramming across C++, Java, and C\#. This specific section must strictly focus on demonstrating and explaining C++ compile-time computation mechanisms.

Requirements:

Task 1 -- Technical Content:

* Provide 2-3 concrete, fully compilable C++ examples of compile-time computation.
* At least one example must demonstrate recursive compile-time computation using traditional C++ templates (e.g., calculating a factorial or Fibonacci sequence).
* At least one example must demonstrate compile-time computation using `constexpr` functions.
* For each example, provide a precise technical explanation of what the compiler actually generates during compilation.
* For each example, explicitly detail the language mechanisms and compiler rules that guarantee the computation executes entirely before runtime.

Task 2 -- Section-Level Framing (STRICT):

* Treat your output strictly as ONE SECTION of an existing, larger chapter.
* Assume the reader already possesses the context of the broader comparison between C++, Java, and C\# generics/templates; do not re-introduce this context.
* Omit any standalone introduction, motivation, or conversational framing (e.g., do not use phrases like "In this section, we will examine..."). Begin immediately with the technical exposition.
* Omit any standalone conclusion, summary of the chapter, or forward-looking remarks.
* Omit citations, bibliographies, or reference lists.

Task 3 -- Content Constraints:

* ENFORCE STRICT SCOPE CONTROL: Address only the requirements stated above. Do not digress into discussions of memory footprint, hardware performance benchmarking, code bloat, garbage collection, or unrelated language comparisons.
* Distinguish clearly between compile-time evaluation (e.g., template instantiation, constant expression evaluation) and runtime execution.
* Use precise terminology (e.g., template instantiation, constant evaluation, monomorphization, substitution) strictly where relevant.
* State compiler or language standard assumptions (e.g., C++11, C++14, C++20) only if strictly required to contextualize the `constexpr` or template mechanisms used in your examples.
* ENFORCE CONCISENESS: The section must be as concise as possible while fully satisfying the technical requirements. Eliminate all padding, restatements of the prompt, and redundant paragraph transitions.

Writing Style Constraint:

Maintain a formal, objective, and precise technical writing style consistent with advanced programming language theory literature. Do not use first-person forms ("I", "we", "our"). Do not include explanations, comments, editorial notes, placeholders, or meta-commentary addressed to the user.

Output Structure:

Output only the raw Markdown content for this specific section. Use appropriate subsection headings (e.g., `###` or `####`) to organize the examples and explanations. Do not include a top-level title for the chapter. Return exactly the text of the section and nothing else.
\end{PromptVerbatim}

\noindent\textbf{Sources:}
\begin{itemize}
\item \textbf{Vandevoorde, Josuttis, Gregor, \emph{C++ Templates: The Complete Guide}} \cite{cpp_templates2017} -- used for template instantiation and recursive compile-time metaprogramming.
\item \textbf{ISO/IEC 14882, C++ Language Specification} \cite{iso14882} -- used as a normative reference for template instantiation.
\item \textbf{cppreference.com, documentation on \texttt{constexpr}} \cite{cppreference_constexpr} -- used for the evaluation of constant expressions at compile time.
\item \textbf{cppreference.com, documentation on templates} \cite{cppreference_templates} -- used for template arguments and the specialization process.
\end{itemize}

% =====================================================================
\section{Prompt 2: SFINAE and Type Traits}
% =====================================================================

\noindent\textbf{Goal}: Generate a detailed academic section of a C++ example using SFINAE that selects a different function overload depending on whether a type is integral vs. floating point vs. a custom class. Also show the equivalent using C++20 concepts. Explain step by step what "substitution failure" means here and why it isn't a compile error.

\noindent\textbf{Model:} ChatGPT-5.5

\noindent\textbf{Prompt}:
\begin{PromptVerbatim}
System Persona:

Act as an expert technical writer and programming language theorist specialized in C++ metaprogramming and compile-time polymorphism.

Goal:

Generate a single, highly focused section for a larger academic and technical chapter on parametric polymorphism and metaprogramming. This section must exclusively demonstrate and explain C++ SFINAE and C++20 Concepts for overload resolution based on type properties.

Requirements:

* Technical Content:
* Provide a working C++ code example using SFINAE (specifically via `std::enable_if`) that successfully selects a different function overload depending on whether a type is an integral type, a floating-point type, or a custom class.
* Provide the equivalent working C++ code example utilizing C++20 Concepts.
* Provide a step-by-step explanation of what "substitution failure" means in this specific context and explicitly detail the mechanism of why it does not result in a hard compile error.

* Content Constraints:
* Distinguish clearly that these are compile-time mechanisms.
* Enforce strict scope control: address only the SFINAE implementation, the C++20 Concepts equivalent, and the explanation of substitution failure. Do not digress into compilation times, memory footprint, code bloat, hardware performance, or other unrelated language comparisons.

Output Structure:

Use standard Markdown code blocks for the C++ code examples.
\end{PromptVerbatim}

\noindent\textbf{Sources:}
\begin{itemize}
    \item \textbf{Stroustrup, \emph{The C++ Programming Language}} \cite{stroustrup2013} -- used for compile-time polymorphism based on templates.
    \item \textbf{cppreference.com} \cite{cppreference} -- used for the behavior of \texttt{std::enable\_if} and C++20 Concepts syntax.
    \item \textbf{ISO/IEC 14882:2020} \cite{iso14882} -- used as a normative reference for SFINAE rules and the C++20 constraint system.
\end{itemize}

% =====================================================================
\section{Prompt 3: Variadic Templates, Template Template Parameters, and CRTP}
% =====================================================================

\noindent\textbf{Goal}: Generate a detailed academic section of other core C++ template metaprogramming techniques: variadic templates, template template parameters, and CRTP (Curiously Recurring Template Pattern). For each, give a short example and explain what class of problems it solves that plain function/class templates cannot.

\noindent\textbf{Model:} ChatGPT-5.5

\noindent\textbf{Prompt}:
\begin{PromptVerbatim}
System Persona:

Act as a technical-academic writer specialized in programming language theory and advanced C++ software engineering.

Goal:

Write a single, highly focused section for a larger chapter on parametric polymorphism and metaprogramming. The section must explain three specific advanced C++ template metaprogramming techniques: variadic templates, template template parameters, and the Curiously Recurring Template Pattern (CRTP).

Requirements:

* Technical Scope: Focus exclusively on variadic templates, template template parameters, and CRTP. Explicitly exclude discussions of SFINAE, type traits, and constexpr, as these are assumed to be covered elsewhere.
* Content Execution: For each of the three covered techniques:
* Provide a short, syntactically correct C++ code example.
* Explain the specific class of problems the technique solves that plain function or class templates cannot resolve.
* Clearly distinguish compile-time mechanisms and state relevant assumptions regarding C++ language/compiler versions (e.g., C++11 for variadic templates) where necessary.

* Strict Scope Control: Do not digress into memory footprint, code bloat, hardware performance benchmarking, or unrelated language comparisons. Address only the mechanisms requested.
\end{PromptVerbatim}

\noindent\textbf{Sources:}
\begin{itemize}
\item \textbf{Vandevoorde, Josuttis, Gregor, \emph{C++ Templates: The Complete Guide}} \cite{cpp_templates2017} -- used for variadic templates, template template parameters, and CRTP.
\item \textbf{cppreference.com} \cite{cppreference} -- used for C++ construct syntax and template specialization.
\item \textbf{Microsoft, C++ templates documentation} \cite{microsoftcpptemplates} -- used as a complementary reference for practical syntax examples.
\end{itemize}

% =====================================================================
\section{Prompt 4: Limits of Java Generics in Replicating C++ SFINAE and Type Traits}
% =====================================================================

\noindent\textbf{Goal}: Generate a detailed academic section on the limits of Java generics in replicating C++ SFINAE and type traits. Show what is achievable (e.g. bounded type parameters like \texttt{<T extends Number>}) and explain precisely why full SFINAE-style compile-time branching on the generic type \texttt{T} is impossible in Java, tying the explanation to type erasure.

\noindent\textbf{Model:} ChatGPT-5.5

\noindent\textbf{Prompt}:
\begin{PromptVerbatim}
System Persona:

Act as an expert in programming language theory and software engineering, specializing in the comparative analysis of parametric polymorphism and metaprogramming across C++, Java, and C\#.

Goal:

Write a single, concise section of an academic report chapter evaluating the limits of Java generics when attempting to replicate C++ SFINAE and type-traits behavior. Specifically, demonstrate the extent of what is achievable in Java (e.g., bounded type parameters like `<T Number extends>`) and explain precisely why full SFINAE-style compile-time branching on a generic type `T` is impossible in Java, explicitly tying this architectural limitation to Java's implementation of type erasure.

Requirements:

* Technical Demonstration: Provide a concise, valid Java code example demonstrating the maximum achievable compile-time constraint mechanism in Java, specifically bounded polymorphism (e.g., `<T Number extends>`).
* Theoretical Explanation: Detail exactly why full compile-time branching on generic types (the equivalent of C++ Substitution Failure Is Not An Error) fails in Java. Connect this limitation mechanically and theoretically to type erasure.
* Terminology: Use precise technical terminology appropriate for programming language theory (e.g., type erasure, bounded polymorphism, monomorphization, compile-time resolution, bounds checking) where directly relevant.
* Scope Control: Address only the replication of SFINAE/type-traits in Java. Do not digress into memory footprint, hardware performance, C\# mechanisms, code bloat, garbage collection, or unrelated language comparisons.
\end{PromptVerbatim}

\noindent\textbf{Sources:}
\begin{itemize}
\item \textbf{Bracha et al., ``Making the Future Safe for the Past''} \cite{bracha1998} -- used for the type erasure model and the limitations of parametric polymorphism in Java.
\item \textbf{Gosling et al., The Java Language Specification} \cite{gosling2014java} -- used for the semantics of the \texttt{extends} constraint and the erasure process.
\item \textbf{Vandevoorde, Josuttis, Gregor, \emph{C++ Templates: The Complete Guide}} \cite{cpp_templates2017} -- used as a comparison for SFINAE substitution and type traits in C++.
\end{itemize}

% =====================================================================
\section{Prompt 5: Limits of C\# Generics in Replicating C++ SFINAE and Type Traits}
% =====================================================================

\noindent\textbf{Goal}: Generate a detailed academic section on the limits of C\# generics in replicating C++ SFINAE and type traits. Show what's achievable using generic constraints (\texttt{where T : ...}) and, if relevant, .NET generic math / static abstract interface members. Explain why C\# gets closer to C++ than Java does, but still can't do arbitrary compile-time computation or true SFINAE, tying it to reification vs. monomorphization.

\noindent\textbf{Model:} ChatGPT-5.5

\noindent\textbf{Prompt}:
\begin{PromptVerbatim}
System Persona:

Act as an expert academic writer and programming language theorist specializing in the comparative analysis of parametric polymorphism and metaprogramming paradigms.

Goal:

Generate a single, concise, and highly technical section for a larger academic chapter. This section must specifically analyze how C\# approximates C++ SFINAE and type-traits using generic constraints and generic math, while explaining its theoretical limitations.

Requirements:

* Technical Execution:
* Replicate the conceptual equivalent of a C++ SFINAE/type-traits example using C\#.
* Demonstrate achievable compile-time/type-check-time constraints using C\# bounded polymorphism (`where T : ...`).
* Illustrate the application of .NET generic math and static abstract interface members to achieve trait-like capabilities.

* Theoretical Analysis:
* Explain why these C\# mechanisms allow it to approach C++ capabilities much more closely than Java does.
* Explain why C\# inherently cannot support true SFINAE (Substitution Failure Is Not An Error) or arbitrary compile-time computation.
* Directly tie this limitation to the architectural differences between .NET's runtime reification of generics and C++'s compile-time monomorphization and AST-level substitution.

* Scope Control: address solely the requested C\# replication of SFINAE/traits, generic constraints, generic math, and the reification vs. monomorphization dichotomy. Do not digress into memory footprint, code bloat, execution performance, garbage collection, or unrelated language features.
\end{PromptVerbatim}

\noindent\textbf{Sources:}
\begin{itemize}
\item \textbf{Microsoft, C\# generics documentation} \cite{dotnetgenerics} -- used for generic constraints and compile-time validation of the \texttt{where} clause.
\item \textbf{Microsoft, .NET Generic Math documentation} \cite{dotnetgenericmath} -- used for static abstract interface members and interfaces such as \texttt{INumber<T>}.
\item \textbf{Vandevoorde, Josuttis, Gregor, \emph{C++ Templates: The Complete Guide}} \cite{cpp_templates2017} -- used as a comparison with SFINAE, type traits, and C++ monomorphization.
\item \textbf{ECMA-335: Common Language Infrastructure} \cite{ecma335_2024} -- used for the CLR execution model and generics reification.
\end{itemize}

% =====================================================================
\section{Prompt 6: Instantiation-Time vs. Declaration-Time Checking}
% =====================================================================

\noindent\textbf{Goal}: Generate a detailed academic section explaining how the difference between C++ templates (checked only at instantiation) and Java/C\# generics (checked at declaration) directly causes specific metaprogramming techniques to be impossible in Java/C\#.

\noindent\textbf{Model:} ChatGPT-5.5

\noindent\textbf{Prompt}:
\begin{PromptVerbatim}
System Persona:

Act as a technical author and programming language theorist writing a highly specialized section for an academic chapter on programming language theory and software engineering.

Goal:

Generate a single, self-contained section of a larger chapter that explains how the difference between C++ templates (checked only at instantiation) and Java/C\# generics (checked at declaration) directly causes specific metaprogramming techniques to be impossible in Java/C\#. You must detail at least one concrete technique that fails specifically due to this difference.

Requirements:

* Focus strictly on the mechanisms of type checking: instantiation-time (C++) versus declaration-time (Java/C\#).
* Explain the direct causal relationship between these checking mechanisms and the limitations placed on metaprogramming in Java and C\#.
* Detail at least one concrete metaprogramming technique (e.g., SFINAE, or structural/duck typing within templates) that is possible in C++ but impossible in Java/C\# because of this specific constraint.
* Include concise, precise code examples to illustrate the technique in C++ and explicitly demonstrate why declaration-time checking prevents its equivalent in Java/C\#.
* Use precise terminology (e.g., instantiation, substitution failure, bounded polymorphism, type erasure, reification) where strictly relevant to the prompt.
* Clearly distinguish between compile-time and runtime mechanisms as they relate to these type-checking models.
\end{PromptVerbatim}

\noindent\textbf{Sources:}
\begin{itemize}
\item \textbf{Vandevoorde, Josuttis, Gregor, \emph{C++ Templates: The Complete Guide}} \cite{cpp_templates2017} -- used for the delayed type checking of templates and SFINAE.
\item \textbf{Bracha, \emph{Generics in the Java Programming Language}} \cite{bracha2004} -- used for the declaration-time checking model of Java generics.
\item \textbf{ECMA-335: Common Language Infrastructure} \cite{ecma335_2024} -- used for the compile-time verification rules of C\# generics.
\end{itemize}

% =====================================================================
\section{Prompt 7: External Metaprogramming Workarounds in Java and C\#}
% =====================================================================

\noindent\textbf{Goal}: Generate a detailed academic section on the external metaprogramming workarounds in Java and C\#. Explain why these are workarounds external to the generics system rather than native metaprogramming.

\noindent\textbf{Model:} ChatGPT-5.5

\noindent\textbf{Prompt}:
\begin{PromptVerbatim}
System Persona:

Act as an academic technical writer specializing in programming language theory and software architecture.

Goal:

Write exactly ONE SECTION of an existing comprehensive chapter on parametric polymorphism and metaprogramming. This specific section must describe the common workarounds used in Java and C\# to simulate the native metaprogramming capabilities of C++ templates, and architecturally explain why these approaches are external to their respective generics systems rather than native metaprogramming.

Requirements:

* Technical Content Specification:
* Detail Java's common workarounds: the use of runtime reflection, compile-time annotation processors, and code generation utilities (e.g., Lombok).
* Detail C\#'s common workarounds: the use of source generators and T4 templates.
* Articulate the structural and theoretical reasons why these mechanisms are considered external workarounds layered over or completely separate from the Java and C\# generics systems, contrasting them with the native, intrinsic compile-time evaluation of C++ templates.
* Clearly differentiate between mechanisms that operate at runtime versus those that operate at compile-time or pre-compilation.

* Content Constraints:
* Use accurate terminology (e.g., compile-time, runtime, reflection, AST manipulation, macro expansion) strictly where relevant.
* Include brief, illustrative code snippets only if strictly necessary to demonstrate a specific structural workaround mechanism; otherwise, rely on rigorous theoretical explanation.
* Strictly enforce scope control: address only the workarounds and their relationship to the generics system. Do not digress into memory footprint, hardware performance benchmarking, garbage collection overhead, or unrelated architectural comparisons.
\end{PromptVerbatim}

\noindent\textbf{Sources:}
\begin{itemize}
\item \textbf{Oracle, Java compiler API documentation} \cite{oracle_java_compiler_apis} -- used for annotation processing and reflection as a runtime mechanism.
\item \textbf{Lombok Project, official documentation} \cite{lombok_project} -- used for code generation via compiler extensions.
\item \textbf{Microsoft, C\# code generation documentation} \cite{microsoft_code_generation} -- used for Roslyn-based Source Generators and T4 Templates.
\item \textbf{Vandevoorde, Josuttis, Gregor, \emph{C++ Templates: The Complete Guide}} \cite{cpp_templates2017} -- used for the template instantiation model as native metaprogramming.
\end{itemize}

% =====================================================================
% BIBLIOGRAPHY
% =====================================================================
\bibliographystyle{plain}
\bibliography{references}

\end{document}